\chapter{Transcatheter Aortic Valve Replacement (TAVR)}

\section*{What Is This Test or Procedure?}
TAVR is a minimally invasive procedure to replace a narrowed aortic valve (aortic stenosis) using a catheter. A collapsible prosthetic valve is guided to the heart and expanded inside the diseased native valve.

\section*{Alternative Names}
TAVR, TAVI, Transcatheter aortic valve insertion

\section*{Principle}
A collapsible prosthetic valve is crimped onto a catheter, advanced to the aortic annulus through the femoral artery or another access route, and expanded by balloon inflation or self-expansion. The new valve pushes the calcified native leaflets aside and takes over valve function.
\section*{History}
Alain Cribier performed the first human TAVR in France in 2002, revolutionizing the treatment of calcific aortic stenosis in non-surgical candidates.

\section*{Why This Test or Procedure Is Ordered}
Ordered for patients with symptomatic severe calcific aortic stenosis who are at intermediate, high, or prohibitive risk for traditional open-heart surgery.

\section*{Is This a Primary or Secondary Test or Procedure}
Primary treatment option for patients with severe aortic stenosis who have high or prohibitive surgical risk.

\section*{How This Test or Procedure Is Performed}
Performed under local anesthesia and sedation (conscious sedation). The catheter is inserted into the femoral artery (transfemoral access) and guided to the aortic valve. Under rapid ventricular pacing, the balloon-expandable or self-expanding prosthetic valve is deployed.

\section*{What Preparation Is Needed}
CT angiography of the chest, abdomen, and pelvis for valve sizing and femoral access planning. Complete blood typing and crossmatch.

\section*{Contraindications and Precautions}
Active endocarditis, inadequate vascular access (femoral arteries too small or calcified), or lack of clinical benefit due to severe comorbidities.

\section*{Risks}
Stroke, vascular access complications, heart block requiring a permanent pacemaker (~10-15\%), or paravalvular leak.

\section*{What Happens After the Test or Procedure}
Recover in the ICU or cardiac step-down unit. Early mobilization is encouraged. Perform an echocardiogram before discharge to evaluate valve function.

\section*{Understanding Your Results}
Echocardiography confirms decrease in transvalvular gradient and improvement in valve area.

\section*{Normal Results}
Low aortic valve mean gradient (<10-15 mmHg); no significant paravalvular leak; normal coronary artery perfusion.

\section*{Follow-up Tests and Procedures}
Echocardiogram in 30 days and then annually; daily low-dose aspirin and/or anticoagulant.

\section*{References}
\begin{enumerate}
  \item MedlinePlus. Transcatheter Aortic Valve Replacement (TAVR). U.S. National Library of Medicine; 2025.
  \item Current Medical Diagnosis and Treatment. McGraw-Hill; 2025.
\end{enumerate}

\clearpage
